\subsection{最多区間の抽出}
% 滞在者の帰宅予測時刻が最多となる30分区間を抽出する．
% これにより多くの人が帰る時間帯を取得する．
昇順に並べられた各滞在者の帰宅予測時刻に対して，幅30分の時間区間を連続的にスライドさせ集計を行い，その中で最も多くのユーザが含まれる時刻の区間（以下，最多区間）を抽出する．
これは，個々の帰宅時刻に着目するのではなく，複数人が帰宅しやすい時間帯を捉える目的で行う．
また，時刻をあらかじめ固定された単位で区切らずに区間を設定すると，ユーザの帰宅時刻のばらつきを許容しつつ，帰宅同行できる可能性が高い時間帯を柔軟に抽出できる．
このようにして得られた最多区間を基に帰宅推奨時刻の妥当性の判定を行い，有効と判断された場合にのみ提示時刻の算出へと進む．


最多区間となる時刻の算出にあたっては，算出結果が複数人の同時帰宅を表す代表時刻として妥当であるかを判定するため，一定の条件を満たす場合のみ有効な結果として扱う．
具体的には，滞在者数が1人以下の場合や，最多区間内に含まれるユーザが1人以下の場合には，帰宅推奨時刻を無効とする．
これは，本システムが複数の滞在者が共に帰宅するきっかけの提供を目的としており，単独行動に近い状況では提示の意義が低いと判断できるためである．
また，このような条件に該当する算出結果については，後続処理や画面表示に用いられないため，データベースには保存せず不要な帰宅推奨時刻のログが蓄積されない設計としている．


区間幅を30分とした理由は，対象環境において，帰宅が想定される時間帯におけるバスの運行間隔が，最も本数の少ない時間帯であっても15分である点にある．
この運行間隔を基準とすると，帰宅時刻をバス1本分前後に調整する範囲に相当し，ユーザにとっての心理的負担は比較的小さいと考えられる．
また，区間幅が広すぎる場合には「一緒に帰る可能性がある集団」としての妥当性が低下し，一方で区間幅が狭すぎる場合には有効と判定される人数が減少し帰宅推奨時刻が成立しにくくなる．
30分という区間幅は，これらの点を踏まえた妥当な設定であると考えられる．
したがって，本研究ではこの区間幅を固定値として採用しているが，利用環境や運用方針によって別の区間幅を用いる余地があると考えられる．


なお，最多区間が複数存在する場合には，より早い時刻の区間を帰宅推奨時刻として採用する．
これは，早い時刻に帰宅予定のユーザは時間の経過とともに帰宅してしまい，以降の提示ではアプローチできなくなる可能性が高いためである．
一方で，遅い時刻に帰宅予定のユーザは再度アプローチできる可能性が残る．
このような特性を踏まえ，帰宅同行の機会損失を最小化する観点から，本システムでは早い時刻の区間を優先的に採用する設計としている．